\scanpage{86}
% archive sheet number: 85
donc
\[
u:\quad s_3\mapsto s''_3\mapsto s'_3\mapsto s_3\tag*{(7)}
\]
Donc on trouve
\[
u^3=\mathrm{id}\qquad\deleted{\unsure{u^2=1}}\quad(u\ne\mathrm{id}\ \text{bien sûr})\tag*{(8)}
\]
Notons que $s_3,s_1,s_2,s'_3$ sont dans un $=$ plan (celui d'équation $x+y+z=1$ !) et ce plan est stable par $\sigma_1$ et $\sigma_2$ (\unsure{idem|donc} par $v=\sigma_1\sigma_2$), \deleted{\unsure{aussi}} donc par \deleted{$\unsure{u}$}\inserted{$v=\sigma_1\sigma_2$}. Un calcul \unsure{connu|comme} montre que
\[
v^5=\mathrm{id}\quad\Longleftrightarrow\quad\boxed{\alpha^2+\alpha-1=0}\tag*{(9)}
\]
Donc si $G$ est le groupe
\[
\begin{aligned}
G&=\left\{\sigma_0,\sigma_1,\sigma_2\ \middle|\ \begin{aligned}&\sigma_0^2=\sigma_1^2=\sigma_2^2=1,\quad\sigma_0\sigma_2=\sigma_2\sigma_1,\\&(\sigma_0\sigma_1)^3=1,\ (\sigma_1\sigma_2)^5=1\end{aligned}\right\}\\
&=\mathrm{Aut}(\text{icosaèdre comb. épinglé}\ldots)
\end{aligned}\tag*{(10)}
\]
alors sous réserve que $\alpha^2+\alpha-1=0$, ce groupe \inserted{fini} $G$ opère \inserted{affinement} sur $\mathbb{E}^3$. Montrons que l'ensemble des transformés de $s_0$, $\mathrm{Dr}(s_0,s_1)$ et $\mathrm{Pl}(s_0,s_1,s_3)$ forme un icosaèdre géométrique ? Cela signifie que sur \unsure{tt|le} corps, les seules relations d'incidence entre ces \deleted{droites} sommets, droites et plans sont celles
